\chapter*{Jak czytać ten dokument}
\addcontentsline{toc}{chapter}{Jak czytać ten dokument}

Kompendium łączy w jednym miejscu trzy dokumenty projektu \textbf{PlayLink Multi-Region}, które
dotąd żyły osobno: dokument wizji, raport z prac na 18 listopada 2026 roku oraz dokumentację techniczną
z repozytorium. Treść została scalona i ujednolicona: każdy temat jest opisany raz, a w innych
miejscach pojawia się tylko odesłanie.

\begin{table}[H]
\centering\small
\renewcommand{\arraystretch}{1.25}
\begin{tabularx}{\textwidth}{@{}p{0.27\textwidth}Y@{}}
\toprule
\textbf{Część} & \textbf{Na jakie pytanie odpowiada} \\
\midrule
I.\ Wizja projektu & Po co powstaje system, dla kogo, w jakim zakresie, przy jakich założeniach, kosztach i ryzykach oraz po czym poznamy sukces. Zawiera słownik pojęć używanych w całym dokumencie. \\
II.\ Stan realizacji & Co zostało faktycznie zbudowane i sprawdzone w filarach F1--F7, jak przebiegały kolejne tygodnie prac i jakie wnioski z tego płyną. \\
III.\ Dokumentacja techniczna & Jak system jest zbudowany: architektura, tożsamość i bezpieczeństwo, model danych, REST API, WebSockety, serwer MCP, frontend, wdrożenie, eksploatacja wieloregionowa, obserwowalność i testy. \\
Załączniki & Uruchomienie lokalne, zasady współpracy w repozytorium i indeks plików źródłowych. \\
\bottomrule
\end{tabularx}
\caption{Układ kompendium.}
\end{table}

Czytelnik zainteresowany celem i oceną projektu przeczyta części I i II; osoba rozwijająca lub
wdrażająca system zacznie od rozdziału o architekturze w części III. Stan opisany w dokumencie
odpowiada gałęzi \code{main} repozytorium z listopada 2026 roku. Fragmenty dokumentacji pierwotnie
napisane po angielsku zostały przetłumaczone; nazwy z kodu (pliki, endpointy, zmienne) pozostawiono
bez zmian.
